\section{Introduction}
\label{sec:intro}

Taking transport, inventory and production decisions jointly improves
supply chain performance \autocite{Stadtler2015SupplyPlanning,
Diaz-Madronero2017ADecisions}. Responsibility for transport planning
usually lies with the supplier
\autocite{Fleischmann2015TransportDistribution}, but in the automotive
industry the manufacturer frequently plans its own inbound transport,
and it does so around an asset material requirements planning does not
represent: the returnable transport item (RTI), the tertiary packaging
in which parts move. The plant returns folded empty RTIs and the
supplier sends them back erected and full, so each flow is the other's
supply and neither can be planned alone. Algorithm-based planning
outperforms a planner's intuition here
\autocite{Diaz-Madronero2014AChain}, most where intuition is weakest:
deciding how many trucks to dispatch when whether a given set of
containers fits inside one is itself a combinatorial question.

That last question is the one the literature does not ask. Transport
capacity in procurement transport planning is a scalar --- freight
volume, or a count of loaded units
\autocite{Hashemi-Amiri2023IntegratedApproach,
Diaz-Madronero2014AChain} --- but containers are not fluid: they stack
in whole layers under a height budget and a crush limit, and their
footprints tile a floor in whole rectangles. A plan drawn against a
volumetric capacity can therefore prescribe loads that no arrangement
of the trailer admits, and the shortfall is discovered on the loading
dock rather than in the model.
Section~\ref{sec:res-abstraction} measures how often, by placeability
proof rather than by argument.

This paper makes four contributions, listed in the order of the
evidence behind them.

\begin{enumerate}
\item A mixed integer formulation, 3D-SCPTOP
  (Section~\ref{sec:model}), joining procurement transport, RTI
  closed-loop and inventory decisions to a full-truckload condition
  stated as a stack-based 3D loading problem, with the per-stack crush
  cap evaluated on the product mix actually loaded in each slot rather
  than on a worst case for its container family.

\item A Dantzig--Wolfe reformulation \autocite{DantzigWolfe1960} and a
  branch-and-price procedure, BP-3D-SCPTOP
  (Sections~\ref{sec:decomp}--\ref{sec:branching}), derived
  structurally rather than by analogy: deleting the quantifier over
  (direction, truck, period) from the block rows \emph{is} the pricing
  problem. The procedure branches over the fleet vector with a
  flow-first relax-and-fix construction as primal heuristic, and its
  lower bounds rest on one proposition: a certified bound of the
  volumetric relaxation is valid for the 3D problem.

\item Accuracy against a reference optimum
  (Section~\ref{sec:res-vol}): on the volumetric problem, where a
  compact formulation solves to proved optimality, the procedure's plan
  and its certified bound bracket that optimum within about a hundredth
  of a percent.

\item A measurement of what a volumetric transport capacity costs
  (Section~\ref{sec:res-abstraction}), with the mechanism behind it and
  a test on the industrial container catalogue that a planner can run
  before solving anything.
\end{enumerate}

The branching hierarchy is finite and exhaustive and every prune is a
proof, so BP-3D-SCPTOP is an exact procedure run under a budget rather
than a heuristic; on the sizes reported here it returns a verified plan
and a certified bound, and proves no optimum.
